% 2026-09-02 typed universal microVM architecture insert.
% Software-semantics theorem with explicit physical boundaries.

\subsection{A typed universal microVM above the W33 routing fabric}

The existing packet VM and recursive HoloBox runtime can be sharpened into a
typed virtual-machine contract by incorporating two corrections from the
current frontier. First, the two order-$51840$ groups used by the architecture
must occupy different instruction namespaces:
\[
 \Sp(4,3)
 \qquad\hbox{and}\qquad
 \PGSp(4,3)\cong W(E_6).
\]
The former is the symplectic central double cover of $\PSp(4,3)$ used for a
Clifford lift; the latter is the projective/similitude outer extension. Equal
orders do not identify the extensions. Second, the two degree-$216$ carriers
over the common $36$-state spread/double-six base are not gauge-equivalent.
They are immutable construction-time machine types:
\[
 \text{circuit-216}\longrightarrow St_{81},
 \qquad
 \text{paired-hemisystem-216}\longrightarrow St_{64}.
\]
Accordingly the reference VM rejects a carrier/module mismatch and exposes no
runtime carrier-conversion opcode.

On top of this typed transport layer, the software-universality core is chosen
deliberately from standard computability theory rather than inferred from the
finite counts. A guest program is an arbitrary finite two-counter machine with
macro instructions
\[
 \operatorname{INC}(r,j),\qquad
 \operatorname{DECJZ}(r,j_{\ne0},j_0),\qquad
 \operatorname{HALT}.
\]
The counters range over $\mathbb N$ in the abstract semantics. Since
two-counter machines are a universal model of computation, this gives
Turing-complete \emph{VM semantics}. Each semantic transition is then wrapped
in a deterministic W33 transport certificate. Reconstructing the symplectic
geometry gives
\[
 \operatorname{SRG}(40,12,2,4),\qquad \operatorname{diam}=2,
\]
so consecutive execution portals require at most two W33 hops. A certificate
contains the pre- and post-state digests, typed instruction, program counters,
immutable carrier, logical-module dimension, route, line buses, and the next
hash-chain root.

The executable witness
\texttt{analysis/w33\_typed\_universal\_microvm.py} verifies the geometry and
the type discipline, checks deterministic replay, and executes a regression
guest
\[
 (c_0,c_1)=(7,11)\longmapsto(18,0)
\]
in $24$ certified steps, with maximum route length two. The companion tests
exercise both carrier types and deliberately attempt the forbidden retyping.

\paragraph{Architectural interpretation.}
The resulting separation is
\[
 \text{guest meaning}
 \;\longrightarrow\;
 \text{typed state transition}
 \;\longrightarrow\;
 \text{W33 placement and route}
 \;\longrightarrow\;
 \text{physical realization}.
\]
Thus routing is proof-carrying realization metadata around a computation, not
a silent redefinition of the guest semantics. The existing content-addressed
recursive runtime supplies the natural scalable state mechanism: a network of
microVMs can be stored as a Merkle-rooted state object and loaded through the
same interface as one leaf VM.

\paragraph{Scope.}
This is a theorem about a software virtual machine whose transitions are
certifiably routed by W33. It does \emph{not} prove that one finite 40-node
physical device contains unbounded memory; an implementation needs recursive or
external scalable storage. It also does not make the finite Clifford kernel
quantum-universal. The photonic quantum architecture still requires the
separately validated non-Clifford resource port, together with physical
threshold, loss, and calibration closure.
